\begin{table}[t]
\centering\footnotesize\setlength{\tabcolsep}{3pt}
\caption{Signed-error certificate at threshold zero. The extrema are exact on the Cartesian Q12 domain; $U_{tr}$ and $U_{te}$ count decisions left uncertified on the fixed training and test inputs. All empirical decision differences are zero. This retrospective analysis does not alter selection.}
\label{tab:signed_certificate}
\begin{tabular}{rrrrrr}
\toprule
$L$ & Bytes & $D_-$ & $D_+$ & $U_{tr}$ & $U_{te}$\\
\midrule
513 & 10,314 & $-10065$ & 8250 & 64 & 12\\
1025 & 20,554 & $-3426$ & 2324 & 26 & 3\\
\bottomrule
\end{tabular}
\end{table}

\paragraph{Signed error interval and threshold agreement.}
The same finite-domain enumeration can retain the signs of the edge errors.
Define $\delta_i(q)=\widetilde f_i(q)-f_i(q)$ and
\begin{equation}
 D_- = \sum_i\min_q\delta_i(q),\qquad
 D_+ = \sum_i\max_q\delta_i(q).
 \label{eq:signed_error_extrema}
\end{equation}
For every admissible input, under the arithmetic assumptions above,
\begin{equation}
 D_-\leq\widetilde Z-Z\leq D_+,
 \qquad Z\in[\widetilde Z-D_+,\widetilde Z-D_-].
 \label{eq:signed_error_interval}
\end{equation}
Indeed, the categorical terms cancel and each numerical error lies between
its enumerated minimum and maximum; summing these inequalities proves the
claim. Because the numerical input domain is a Cartesian product, the lower
and upper error endpoints are attained by choosing an extremizing input
separately for each edge. Thus these are exact global extrema on that
Cartesian domain, even if some combinations do not correspond to physical
flows. Their computation costs ten enumerations of 8,193 inputs, not an
enumeration of $8193^{10}$ combinations.

For the decision rule $\mathbf1[Z\geq\tau]$, agreement is certified by either
of the following conditions, with equality assigned to the positive side:
\begin{align}
 &\widetilde Z\geq\tau\quad\text{and}\quad\widetilde Z-D_+\geq\tau,
 \quad\text{or}\nonumber\\
 &\widetilde Z<\tau\quad\text{and}\quad\widetilde Z-D_-<\tau.
 \label{eq:signed_threshold_certificate}
\end{align}
For the tables evaluated here, $D_-\leq0\leq D_+$, so the interval contains
$\widetilde Z$ and the condition reduces to checking its endpoints.
The same argument preserves a fixed risk tier if $\widetilde Z$ and both
endpoints lie in the same half-open threshold interval. This is a direct corollary of
the error interval, not an evaluation of a risk-tier application.

Table~\ref{tab:signed_certificate} summarizes the signed bounds for both grids.
For the frozen $L=1025$ arrays, $(D_-,D_+)=(-3426,2324)$.
The interval based on the LUT score certifies 168,808 of 168,834 training
decisions and 42,206 of 42,209 test decisions. The remaining 26 and 3 rows,
respectively, agree empirically but are not certified by this interval test.
The tighter analysis leaves the selected table and all measured firmware
unchanged. Certificate evaluation is performed on the host; no runtime
guard or fallback has been added to the timed device kernels. Its enumerated edge errors also agree with 81,930 independent
one-coordinate probes of the host-compiled C implementation.

For a finite labeled cohort of $N$ rows, let $U$ be its uncertified count.
Since predictions can differ only on those rows,
$|\mathrm{Acc}_{\rm LUT}-\mathrm{Acc}_{\rm coeff}|\leq U/N$.
For balanced accuracy, with $N_y$ rows and $U_y$ uncertified rows of class $y$,
\begin{equation}
 |\mathrm{BA}_{\rm LUT}-\mathrm{BA}_{\rm coeff}|
 \leq \tfrac12\left(U_0/N_0+U_1/N_1\right).
 \label{eq:certificate_metric_bound}
\end{equation}
For the 42,209 test rows, $(N_0,N_1)=(10000,32209)$ and the signed-interval
test leaves $(U_0,U_1)=(0,3)$. Its worst-case absolute accuracy and balanced
accuracy changes are therefore at most 0.00711 and 0.00466 percentage points,
respectively. These bounds compare the two fixed integer representations
on this cohort; they do not guarantee correctness against ground truth or
performance on an unseen distribution. The observed changes are zero.

\paragraph{Offline use in export verification.}
The certificate can serve as an export check before a firmware release.
A release record can bind the coefficient and LUT hashes, feature encoding,
arithmetic rules and decision threshold to the enumerated error interval.
The interval then defines which LUT scores have certified agreement with
the reference. On a declared audit cohort, the same check reports the
uncertified rows for explicit coefficient comparison. This separates a
whole-domain error bound from agreement verified only on observed inputs.
A future automated build could reject an export when its uncertified count
exceeds a budget fixed on training or validation inputs before test access.
That is a possible use of the analysis, not an implemented CI gate or the
selection rule used in this study. A cohort budget does not bound the
uncertified frequency on a new network. For online handling of uncertain
scores, a guard and a specified fallback would need separate implementation
and memory/latency measurement. The present offline check adds no such
protection to the measured firmware.

\paragraph{Retrospective memory--uncertainty comparison.}
The $L=513$ table occupies 10,314 bytes and also has zero observed training
and test decision differences. Its signed error interval is $[-10065,8250]$,
and the LUT-score interval leaves 64 training and 12 test decisions
uncertified. Increasing to $L=1025$ adds 10,240 parameter bytes (a factor of
1.993) and reduces these counts to 26 and 3. Under the original symmetric
reference-margin condition, the corresponding counts are 95 to 34 on
training and 21 to 4 on test. For these two fixed tables, the extra memory narrows the uncertainty
region; it does not purchase a complete decision-preservation guarantee
or establish the minimum storage required for an uncertainty budget.
The $L=513$ test analysis is retrospective and is not used to change the
training-only selection or the frozen deployment.
